% Pursuing Stacks (cotes 134-2 to 134-8): archive cote 134-3, site page 81 (PDF page 82), typescript page 81.
% typescript. Draft compared: ps2.tex:4210-4279.
% AI-transcribed by Claude Opus 5.5 (claude-opus-5-5, adaptive thinking, effort high), Message Batches API, batch pilot (msgbatch_0189Nhf1xB615YWnztubtLfc), 2026-11-12; not yet checked by eye.
\scanpage{134-3/81}
\mspage{81}

if we grant the ``insp\xout{i}iring assumption'',\xout{xx} implying that any autoequivalen\xout{x}
ce of(Hot)  is isomorphic to the identity functor  -  in this case it is
clear that even the weaker form (i')  of (i),  demanding only that
$i_Aj_A$ \inserted{induce} \deleted{be} an \inserted{auto}equivalence \inserted{of (Hot)} (and \inserted[typed]{nothing on} \xout{noxxthxx} either factor $i_A$ , $j_A$ \inserted{in (5)}),implies (v).
Also,  When we assume moreover A aspheric,\xout{xxxxxxx} it is clear that (i)
(and even (i')) implies (vii) i.e. all other conditions.  Thus,  instead \xout{x}
of the assumption  on (Hot),  it would be enough to know that the particu-%
lar autoequivalence \xout{$i_Aj_A$} of (Hot) induced by $i_Aj_A$  transforms the final
element of (Hot)  (respresented by the element e of (Cat)) into an isomor-%
phic one,\inserted{\unsure{(98)}} which looks like a very slight strenghtening of (i) indeed.
\margin{\hand{(98)} * 27.3. But this is true for any equivalence of categories -- I'm really being very dull!}

In any case,  the conditions (ii) to (viii) are equivalent,  and
equivalent to  (i)  (or (i'))  \uline{plus}\inserted{\unsure{(99)}} asphericity of A.  The formally
\margin{\hand{(99)} ** The ``plus'' is unnecessary visibly, see note above.}
strongest form is (iii),  the formally weakest one (with the exception
of (i) or (i')\hand{)} \deleted{plus \hand{(**)} asphericity of A)} is (v),  which is also the one
which looks the most concrete,  namely amenable to practical verification.
This was indeed what at the very beginning was attractive in this conditi-%
on,  in comparison to the first one that came to my mind when introducing
the notion of a modelizer and of model preserving functors  -  namely \xout{that}
merely that  $i_A$  should induce an equivalence between the localized cate-%
gories,  or equivalently,  that A\hand{$\hat{}$}should be a ``modelizer'' and $i_A$ should
be model-preserving.  That we should be more demanding and ask for $j_A$ to
be equally model preserving crept in first rather timidly  -  and I still
dont know (and did\unsure{nt|n't} really stop to think)  if \inserted{the first} one implies the other\deleted{\unsure{(***)}}\hand{(50)}.
\margin{(***) In fact, it is the condition that $j_A = i_A^*$ \deleted{\illegible{1}}, \unsure{not $i_A$}, should be model preserving, which is ``the right'' condition on A, equivalent to A being a weak test category, \deleted{\illegible{3}} if A is assumed aspheric.}

In any case,  as far as checking a property goes,  I would consider
(v) to be \uline{the} handy definition of a weak test category,  whereas (iii)
is the best,  when it goes to making use of the fact that A i\xout{s}s indeed a
weak test category.  As for (i),  \deleted{(supplemented by asphericity to be safe)}
it corresponds to the main intuitive content of the notion,  which means
that homotopy types can be``modelled'' by elements of A\^{} \xout{xxxxxxxx} ,  using
$i_A$ for describing which homotopy type is describe\xout{e}d by an element F of A\^{},
and using $j_A$ for getting a model in A\hand{$\hat{}$}for a given homotopy type, descri-%
be\xout{e}d by an object C in (Cat).

b)  \inserted{\deleted{\unsure{Weak}}} \uline{Test categories}\inserted{and local test categories}.  Let A still be a small category.  T\unsure{\over{h}{h}}en the follo-%
wing conditions are equivalent:

\deleted{i)  A is a weak test category,  and so are all induced categories}
\deleted{$A_{/a}$  with a in A.}
\deleted{ii)  A is aspheric,  and for any aspheric element F in A\^{}, $A_{/F}$ is}
\deleted{a weak test category.}
